% The three required disclosures (about 150 words each), shared by the pre-analysis plan and the final manuscript.
% \finalonly{...} text appears only in the final manuscript, which defines it before this file.
\providecommand{\finalonly}[1]{}
\section*{Required statements}

\subsection*{Integrity statement}
Plagiarism of text: every source we draw on is cited, nothing is copied without quotation marks, and references were checked against OpenAlex. Plagiarism of ideas: borrowed methods are credited where they appear (BALD, DOSE, cognitive-model pretraining), and all code is our own. Bias: we fixed hypotheses and decision rules before running experiments, list every deviation, and report each dataset separately because the samples are Swiss, German and US adults. Inaccuracies: the public Basel--Berlin file codes 0--10 ratings as 1--11, which we correct; Basel--Berlin informed our design, so its results are exploratory, and Twin-2K-500 was analyzed once, as the confirmatory test.\finalonly{ Selective reporting: results that went against our hypotheses, H1 and H2, are reported as fully as those that supported them, and every number, figure and table in the paper is produced by the submitted code from the public data.}

\subsection*{AI use reflection}
We used Claude Opus 5.5 to search the literature, draft and debug analysis code, and edit our writing. We used it because it made routine work such as data loading and plotting much faster, leaving time to compare four model classes. The research question, model classes, hypotheses and decision rules are ours, and we reviewed all code and text before using it. The tool also made confident errors, so every citation was checked against OpenAlex and every number comes from code we ran. For our learning, checking its output pushed us to work through expected information gain and identification ourselves; the risk is accepting code we do not understand, which we guard against by reading and testing each function. In the future we would ask AI tools to critique a design before we build it, and the course could ask students to submit their AI transcripts with their code.

\subsection*{Impact statement}
Fewer, better-chosen questions could make risk profiling quicker and more accurate, so clients who would skip a long questionnaire still get advice that fits them. The methods could be misused: an adaptive questionnaire could be steered to label clients more risk tolerant than they are, or to profile them for marketing. Mitigations are to publish the question-selection rule and the model, keep the self-rating as a cross-check, and let clients see and correct their profile; our code is submitted for audit. Lottery choices are only a proxy for portfolio decisions, our samples are Swiss, German and American, and adaptive questions made predictions overconfident on one dataset, so the findings are not ready for use with real clients. We collect no new data: all datasets are public and de-identified, used under their licenses and cited, as Stanford's integrity standards require. Everything runs on a laptop, so the environmental cost is negligible.

